\chapter{Total Hip Replacement}

\section*{What Is This Test or Procedure?}
Total hip replacement removes damaged portions of the hip joint and replaces them with artificial components to reduce pain and restore function.

\section*{Why This Test or Procedure Is Ordered}
It is commonly performed for severe osteoarthritis, inflammatory arthritis, fracture, avascular necrosis, or other joint destruction causing persistent pain and disability.

\section*{How This Test or Procedure Is Performed}
Under regional or general anesthesia, the damaged femoral head is replaced with a stem and ball, and the damaged socket is resurfaced with an artificial cup and liner. Different surgical approaches may be used.

\section*{Risks and Recovery}
Infection, blood clots, dislocation, fracture, nerve injury, leg-length difference, loosening, and revision surgery are possible. Early mobilization, precautions, and rehabilitation support recovery.

\section*{References}
\begin{enumerate}
  \item MedlinePlus. Hip Replacement. U.S. National Library of Medicine; 2025.
  \item Current Medical Diagnosis and Treatment. McGraw-Hill; 2025.
\end{enumerate}
\clearpage
